% ══════════════════════════════════════════════════════════════════
% INTRODUCTION GÉNÉRALE
% ══════════════════════════════════════════════════════════════════
\newpage
\pagenumbering{arabic}
\setcounter{page}{1}
\chapter*{Introduction Générale}
\addcontentsline{toc}{chapter}{Introduction Générale}

L'informatique a bel et bien révolutionné le monde actuel où nous vivons, particulièrement à partir des années quatre-vingt-dix. Elle est présente partout : dans le travail, les communications, les transports, le commerce, les services, etc. On est devenu presque dépendant de cet outil qui a pratiquement envahit notre existence. Certes, l'informatique a beaucoup contribué à l'amélioration des conditions de notre vie. Elle nous a permis d'exécuter rapidement et efficacement des tâches, autrefois très lentes et compliquées. C'est pourquoi, elle est intégrée dans quasiment tous les secteurs d'activité : l'administration, l'industrie, le commerce, la finance, etc.

Dans le cadre de ma formation à l'université EMSI en génie Informatique, j'ai effectué un stage de fin d'études de 6 mois chez HPS.

HPS est un leader mondialement reconnu dans l'édition de solutions de paiement électronique dédiées aux institutions financières, elle fournit des solutions et des services à haute valeur ajoutée autour de sa suite de solutions de paiement électronique PowerCARD.

Ce stage a pour but de valider mes compétences et savoir acquis durant ma formation en accompagnant l'équipe de test d'HPS dans leurs activités de test, d'automatisation de test.

Cette expérience m'a permis de mettre en pratique tout ce que j'ai appris durant ma formation et aussi d'avoir une vision globale sur les activités de l'équipe de test de HPS.

Le rapport présente en premier lieu l'entreprise, le service d'accueil et l'équipe avec laquelle j'ai travaillé, puis le contexte global des missions réalisées ainsi que les détails de chacune d'elle en expliquant le contexte, problématique, la méthodologie utilisée et les résultats obtenus pour enfin finir avec une conclusion.
